\documentclass[11pt]{amsart}

\usepackage[T1]{fontenc}
\usepackage{lmodern}
\usepackage{microtype}
\usepackage{amsmath,amssymb,amsthm}
\usepackage{booktabs}
\usepackage[hidelinks]{hyperref}

\hypersetup{
  pdftitle={The Minimum Dimension of a Maximal Face of the 6x6 Completely Positive Cone}
}

\newtheorem{theorem}{Theorem}
\newtheorem{lemma}[theorem]{Lemma}
\theoremstyle{remark}
\newtheorem{remark}[theorem]{Remark}

\newcommand{\CP}{\mathcal{CP}}
\newcommand{\COP}{\mathcal{COP}}
\newcommand{\Smat}{\mathcal{S}}
\newcommand{\R}{\mathbb{R}}
\newcommand{\supp}{\operatorname{supp}}
\newcommand{\low}{\operatorname{low}}
\newcommand{\cone}{\operatorname{cone}}
\newcommand{\ip}[2]{\langle #1,#2\rangle}
\newcommand{\doi}[1]{\href{https://doi.org/#1}{doi:#1}}

\title[The minimum dimension in $\CP^6$]
{The Minimum Dimension of a Maximal Face of the
$6\times6$ Completely Positive Cone}

\date{}

\subjclass[2020]{15B48, 52A20, 90C26}
\keywords{Completely positive cone, copositive cone, maximal face,
minimal zero, extreme ray}

\begin{document}

\begin{abstract}
Let $\low(n)$ denote the minimum dimension of a maximal proper face of the
$n\times n$ completely positive cone. We prove
\[
\low(6)=7.
\]
Holmgren and Zhang established the upper bound by exhibiting a
$7$-dimensional maximal face. The contribution here is the matching global
lower bound, obtained from the classification of the extreme rays of
$\COP^6$ and a linear-independence argument for rank-one matrices generated
by zeros. This resolves the $n=6$ instance of the even-dimensional open
problem stated by Kostyukova and Tchemisova.
\end{abstract}

\maketitle

\section{Preliminaries and statement}

Let $\Smat^n$ be the space of real symmetric $n\times n$ matrices, equipped
with $\ip{A}{X}=\operatorname{tr}(AX)$. The completely positive and
copositive cones are
\begin{gather*}
\CP^n=\cone\{xx^{\mathsf T}:x\in\R_+^n\},\\
\COP^n=(\CP^n)^*
 =\{A\in\Smat^n:x^{\mathsf T}Ax\ge0\ \text{for all }x\in\R_+^n\}.
\end{gather*}
For a maximal proper face $F$ of $\CP^n$, its dimension is the dimension of
its linear span, and
\[
\low(n)=\min\{\dim F:F\text{ is a maximal proper face of }\CP^n\}.
\]

\begin{theorem}\label{thm:main}
Every maximal proper face of $\CP^6$ has dimension at least $7$, and
$\CP^6$ has a maximal face of dimension $7$. Consequently,
\[
\low(6)=7.
\]
\end{theorem}

The problem of determining the exact value of the tight lower bound on the
dimension of a maximal face of $\CP^n$ was raised by Dickinson
\cite{Dickinson}. Holmgren and Zhang exhibited a $7$-dimensional maximal
face of $\CP^6$ from an exposed Case~19 ray of $\COP^6$ and obtained the
formal bound $6\le\low(6)\le7$
\cite[Lemma~4.7 and Theorems~4.8--4.9]{HZ}. Kostyukova and Tchemisova subsequently proved $n\le\low(n)\le n+3$ for even
$n\ge6$ \cite{KT}, and left the exact value of $\low(n)$ open in that range
\cite[Section~5]{KT}; at $n=6$ their estimate reads $6\le\low(6)\le9$ and so
does not improve on the bound above. The interval that
Theorem~\ref{thm:main} closes is therefore $\{6,7\}$, and the new assertion
is the global lower bound $\low(6)\ge7$.

For $A\in\COP^n\setminus\{0\}$, put
\[
F_A=\CP^n\cap A^\perp.
\]
If $X=\sum_r x_rx_r^{\mathsf T}\in\CP^n$, then
$\ip{A}{X}=\sum_r x_r^{\mathsf T}Ax_r$, and every summand is nonnegative.
Hence
\begin{equation}\label{eq:facezeros}
F_A=\cone\{xx^{\mathsf T}:x\in\R_+^n\setminus\{0\},\ x^{\mathsf T}Ax=0\}.
\end{equation}

\begin{lemma}[Extreme exposer]\label{lem:exposer}
Let $K$ be a proper cone in a finite-dimensional inner-product space.
Every maximal proper face $F$ of $K$ has the form
\[
F=K\cap y^\perp
\]
for some $y$ generating an extreme ray of $K^*$.
\end{lemma}

\begin{proof}
Let $x$ lie in the relative interior of $F$. If $z\in K^*$ satisfies
$\ip{z}{x}=0$, then for each $w\in F$ the point $x-\varepsilon(w-x)$ lies in
$F$ for all small $\varepsilon>0$; since $z$ is nonnegative on $F$ and
vanishes at $x$, this forces $\ip{z}{w}=0$. Hence the conjugate face
\[
F^\triangle=K^*\cap F^\perp=K^*\cap x^\perp
\]
is the set on which the linear functional $\ip{\cdot}{x}$, nonnegative on
$K^*$ because $x\in K$, vanishes; so $F^\triangle$ is an exposed face of
$K^*$. It is nonzero, since $F$ is a proper face and therefore $x$ lies on
the boundary of $K$, so a hyperplane supporting $K$ at $x$ supplies a
nonzero element of $K^*\cap x^\perp$.

As $K^*$ is again a proper cone, $F^\triangle$ is pointed and hence has an
extreme ray $\R_+y$. Since $F^\triangle$ is a face of
$K^*$, the ray $\R_+y$ is extreme in $K^*$. Moreover,
\[
F\subseteq K\cap y^\perp\subsetneq K,
\]
so maximality gives $F=K\cap y^\perp$.
\end{proof}

A nonzero vector $u\in\R_+^n$ is a \emph{zero} of $A\in\COP^n$ if
$u^{\mathsf T}Au=0$. It is \emph{minimal} if no zero has support strictly
contained in $\supp u$. We normalize zeros by $\mathbf1^{\mathsf T}u=1$.

\begin{lemma}[Zero-dimension bound]\label{lem:zeros}
Let $u_1,\dots,u_p$ be distinct normalized zeros of $A\in\COP^n$ that span
$\R^n$. Then $\dim F_A\ge n+1$ if either
\begin{enumerate}
\item $p\ge n+1$, or
\item $p=n$ and $u_i^{\mathsf T}Au_j=0$ for some $i\ne j$.
\end{enumerate}
\end{lemma}

\begin{proof}
After relabeling, let $u_1,\dots,u_n$ be linearly independent and put
$U=(u_1\ \cdots\ u_n)$. The invertible linear map
\[
\Phi(X)=U^{-1}XU^{-\mathsf T}
\]
sends $u_i u_i^{\mathsf T}$ to $e_i e_i^{\mathsf T}$, so these $n$
matrices are linearly independent.

If $p\ge n+1$, write $u_{n+1}=Uc$. Since the zeros are distinct and
normalized, $c$ has at least two nonzero coordinates. Thus
$\Phi(u_{n+1}u_{n+1}^{\mathsf T})=cc^{\mathsf T}$ has a nonzero
off-diagonal entry and is independent of the $n$ diagonal matrices.

If $p=n$ and $u_i^{\mathsf T}Au_j=0$, then $u_i+u_j$ is a zero and
\[
\Phi\bigl((u_i+u_j)(u_i+u_j)^{\mathsf T}\bigr)
=(e_i+e_j)(e_i+e_j)^{\mathsf T},
\]
which again supplies an independent off-diagonal direction. Equation
\eqref{eq:facezeros} completes the proof.
\end{proof}

\section{Proof of the main theorem}

\begin{proof}[Proof of Theorem~\ref{thm:main}]
Let $F$ be a maximal proper face of $\CP^6$. By
Lemma~\ref{lem:exposer}, $F=F_A$ for a matrix $A$ generating an extreme ray
of $\COP^6$.

A copositive matrix is called exceptional if it is not the sum of a
positive-semidefinite matrix and an entrywise nonnegative matrix. Afonin,
Hildebrand, and Dickinson classified all extreme rays of $\COP^6$
\cite[Theorem~5.1]{AHD}. Up to permutation and positive diagonal
congruence, such a ray is nonexceptional, is obtained by padding an
exceptional extreme ray of $\COP^5$ with a zero row and column, or belongs
to one of Cases~1--19. These congruences induce automorphisms of $\CP^6$
and preserve $\dim F_A$.

\smallskip
\noindent\emph{Nonexceptional rays.}
They are generated by $E_{ii}=e_ie_i^{\mathsf T}$, by
$E_{ij}=e_ie_j^{\mathsf T}+e_je_i^{\mathsf T}$ for $i<j$, and by
$aa^{\mathsf T}$ with $a$ having both positive and negative coordinates
\cite[Section~2]{AHD}. Their zero sets give
\[
\dim F_{E_{ii}}=15,
\qquad
\dim F_{E_{ij}}=15+15-10=20,
\qquad
\dim F_{aa^{\mathsf T}}=15.
\]
Indeed, the first face spans the symmetric matrices supported outside row
and column $i$. For $E_{ij}$, the zero set is the union of the coordinate
hyperplanes $x_i=0$ and $x_j=0$; the corresponding $15$-dimensional
symmetric subspaces intersect in a $10$-dimensional subspace. Finally, the
zeros of $aa^{\mathsf T}$ form $a^\perp\cap\R_+^6$, which has nonempty
relative interior in the $5$-dimensional space $a^\perp$, and their
rank-one squares span all symmetric forms on that space.

\smallskip
\noindent\emph{Exceptional rays.}
An exceptional extreme matrix is irreducible with respect to the
positive-semidefinite cone. Otherwise, for some $0\ne P\succeq0$ and
$\varepsilon>0$, the decomposition
$A=(A-\varepsilon P)+\varepsilon P$ inside $\COP^6$ would force $A$ to be
positive semidefinite by extremality. Hildebrand's criterion therefore
implies that the minimal zeros of $A$ span $\R^6$
\cite[Theorem~4.5]{Hildebrand}.

Suppose first that $A=B\oplus0$ is a zero-padded exceptional extreme matrix
of order $5$; this is component~O5 of \cite[Table~1]{HA}, and it is the one
component we argue directly rather than read off that table.\footnote{With
$I_\alpha,J_\alpha$ as defined below, the O5 row records
$J_\alpha\setminus I_\alpha=\emptyset$ for $\alpha\le5$. Since the sixth row
of $A$ vanishes, $(Au_\alpha)_6=0$ for every $\alpha$, so $6\in J_\alpha$
always and the row understates $J_\alpha$ for $\alpha\le5$. The argument
given here uses only $Ae_6=0$ and is independent of that row.}
Then $e_6$ is a minimal zero, and every other minimal zero has
sixth coordinate zero: deleting a positive sixth coordinate would otherwise
produce a zero with strictly smaller support. Thus every other minimal zero
$u$ satisfies $u^{\mathsf T}Ae_6=0$. Since the minimal zeros span $\R^6$,
there are at least six. If there are at least seven,
Lemma~\ref{lem:zeros}(1) applies; if there are exactly six,
Lemma~\ref{lem:zeros}(2) applies to $e_6$ and any other minimal zero. Hence
$\dim F_A\ge7$.

It remains to treat Cases~1--19. Let $u_1,\dots,u_p$ be the normalized
minimal zeros and define
\[
I_\alpha=\supp u_\alpha,
\qquad
J_\alpha=\{k:(Au_\alpha)_k=0\}.
\]
For any zero $u$ of a copositive matrix, $Au\ge0$: applying copositivity to
$u+t e_k$ and differentiating at $t=0^+$ gives $(Au)_k\ge0$. Consequently,
\begin{equation}\label{eq:criterion}
u_\beta^{\mathsf T}Au_\alpha=0
\quad\Longleftrightarrow\quad
I_\beta\subseteq J_\alpha.
\end{equation}
Since $u_\alpha^{\mathsf T}Au_\alpha=0$ and $Au_\alpha\ge0$, every $k\in
I_\alpha$ has $(Au_\alpha)_k=0$; thus $I_\alpha\subseteq J_\alpha$.

The extended minimal-zero support data of Hildebrand and Afonin
\cite[Table~1]{HA} are reproduced in Table~\ref{tab:audit}, in their
labelling both of the components and of the minimal zeros within each
component. Their table lists $22$ main components: O5 is the zero-padded
case treated above, and the other $21$ are those below. Only $I_\alpha$ and
$J_\alpha\setminus I_\alpha$ are recorded, since $I_\alpha\subseteq J_\alpha$
gives $J_\alpha=I_\alpha\cup(J_\alpha\setminus I_\alpha)$; a subset of
$\{1,\dots,6\}$ is written as the string of its elements, so $146$ denotes
$\{1,4,6\}$. For each component with $p=6$ the last column names a pair
$\beta\ne\alpha$ with $I_\beta\subseteq J_\alpha$, which by
\eqref{eq:criterion} is the hypothesis of Lemma~\ref{lem:zeros}(2).

\begin{table}[ht]
\centering
\scriptsize
\setlength{\tabcolsep}{4pt}
\begin{tabular}{@{}ccllc@{}}
\toprule
Component & $p$ & $I_\alpha$ & $J_\alpha\setminus I_\alpha$ & certificate \\
\midrule
$1$ & $6$ & $12,13,14,25,36,456$ & $345,246,23,16,15,\emptyset$ & $I_{2}\subseteq J_{1}$ \\
$2$ & $6$ & $12,13,14,25,356,456$ & $345,246,23,16,\emptyset,\emptyset$ & $I_{2}\subseteq J_{1}$ \\
$3$ & $6$ & $12,13,14,256,356,456$ & $345,24,236,\emptyset,\emptyset,\emptyset$ & $I_{2}\subseteq J_{1}$ \\
$4$ & $6$ & $12,13,24,345,156,456$ & $346,26,15,\emptyset,\emptyset,\emptyset$ & $I_{2}\subseteq J_{1}$ \\
$5$ & $6$ & $12,13,145,246,346,456$ & $35,256,\emptyset,\emptyset,\emptyset,\emptyset$ & $I_{2}\subseteq J_{1}$ \\
$6$ & $6$ & $12,13,245,345,246,356$ & $34,256,\emptyset,\emptyset,\emptyset,1$ & $I_{2}\subseteq J_{1}$ \\
$7$ & $6$ & $15,26,123,234,345,456$ & $24,13,6,\emptyset,\emptyset,\emptyset$ & $I_{3}\subseteq J_{2}$ \\
$8$ & $6$ & $12,134,135,246,346,256$ & $36,5,4,\emptyset,\emptyset,\emptyset$ & $I_{3}\subseteq J_{2}$ \\
$9.1$ & $6$ & $12,134,135,246,346,456$ & $36,\emptyset,\emptyset,5,\emptyset,2$ & $I_{6}\subseteq J_{4}$ \\
$9.2$ & $6$ & $12,134,135,246,346,456$ & $356,\emptyset,2,\emptyset,\emptyset,\emptyset$ & $I_{3}\subseteq J_{1}$ \\
$10$ & $6$ & $12,134,135,246,356,456$ & $36,\emptyset,\emptyset,5,\emptyset,2$ & $I_{6}\subseteq J_{4}$ \\
$11$ & $6$ & $123,124,125,136,246,346$ & $5,5,34,\emptyset,\emptyset,5$ & $I_{3}\subseteq J_{2}$ \\
$12$ & $6$ & $123,124,125,136,246,356$ & $\emptyset,5,4,\emptyset,\emptyset,4$ & $I_{3}\subseteq J_{2}$ \\
$13.1$ & $6$ & $123,234,345,456,156,126$ & $4,15,26,3,\emptyset,\emptyset$ & $I_{2}\subseteq J_{1}$ \\
$13.2$ & $6$ & $123,234,345,456,156,126$ & $4,1,6,3,2,5$ & $I_{2}\subseteq J_{1}$ \\
$14$ & $7$ & $12,13,14,25,45,36,56$ & $345,246,235,146,126,15,234$ & --- \\
$15$ & $7$ & $12,134,135,146,256,356,456$ & $34,2,\emptyset,\emptyset,\emptyset,\emptyset,\emptyset$ & --- \\
$16$ & $7$ & $123,124,125,136,246,346,356$ & $\emptyset,\emptyset,\emptyset,\emptyset,\emptyset,5,4$ & --- \\
$17$ & $7$ & $123,124,125,136,246,356,456$ & $\emptyset,\emptyset,\emptyset,\emptyset,\emptyset,4,3$ & --- \\
$18$ & $8$ & $123,234,345,145,125,346,146,126$ & $\emptyset,\emptyset,6,6,6,5,5,5$ & --- \\
$19$ & $6$ & $345,145,125,123,156,2346$ & $\emptyset,6,\emptyset,\emptyset,4,\emptyset$ & $I_{5}\subseteq J_{2}$ \\
\bottomrule
\end{tabular}
\caption{The classification data needed for the lower bound, from
\cite[Table~1]{HA}.}
\label{tab:audit}
\end{table}

Thus Cases~14--18 satisfy Lemma~\ref{lem:zeros}(1), while every other
component satisfies Lemma~\ref{lem:zeros}(2) by \eqref{eq:criterion}. It
remains to check that this covers every exceptional extreme matrix, and not
only the generic ones. Two features of the classification, both recorded in
\cite[Section~2]{HA}, make this immediate.

First, the extreme matrices of $\COP^6$ are classified in \cite{AHD} ``with
respect to their minimal zero support sets''. The classification therefore
refines the partition by minimal zero support set, so the collection
$\{I_1,\dots,I_p\}$, and in particular $p$, is constant on each piece (two
pieces may share one collection, as with 9.1/9.2 and 13.1/13.2). In particular the
values $p=7$ for Cases~14--17 and $p=8$ for Case~18 hold throughout those
pieces, and Lemma~\ref{lem:zeros}(1) applies to all of their matrices without
any certificate.

Second, the sets $J_\alpha$ can only grow away from the generic locus. The
matrices of a piece are given by analytic expressions in a scaling matrix and
some angles $\phi_j$ \cite[Section~5]{AHD}, and whether $(Au_\alpha)_i=0$
depends only on the $\phi_j$; ``larger complementary supports correspond to
additional equality relations on $\phi_j$, and the corresponding component of
extreme matrices has lower dimension''. Moreover ``almost all matrices in
each piece belong to a single component, while the rest is located on the
boundary of this component''. The component tabulated in
Table~\ref{tab:audit} is thus the generic one, and since $(Au_\alpha)_i$ is
continuous in the parameters, an entry vanishing on the generic locus vanishes
on its closure, that is, on the whole piece. Hence at every matrix of the
piece each $J_\alpha$ contains the set tabulated for it, and the certified
inclusion $I_\beta\subseteq J_\alpha$ holds a fortiori. Cases~$9$ and~$13$ are
the two that decompose into two sub-cases each, listed as $9.1,9.2$ and
$13.1,13.2$ in \cite[Section~3.4]{AHD}: these share a minimal zero support
set but are parametrized separately, and each of the four is a piece with its
own main component. All four appear in Table~\ref{tab:audit} with a
certificate. Therefore
every exceptional extreme ray also satisfies
\[
\dim F_A\ge7.
\]
We conclude that $\low(6)\ge7$.

Finally, Holmgren and Zhang proved that a Case~19 matrix generates an
exposed ray of $\COP^6$ whose conjugate maximal face has dimension $7$
\cite[Lemma~4.7 and Theorem~4.8]{HZ}. Hence $\low(6)\le7$, and equality
follows.
\end{proof}

\begin{remark}
The proof uses the full extreme-ray classification only for the lower bound
and one exposed Case~19 ray for the upper bound. It makes no exposedness
claim for the other exceptional rays. The exact values of $\low(n)$ for
even $n\ge8$ remain open.
\end{remark}

\begin{thebibliography}{9}

\bibitem{AHD}
A.~Afonin, R.~Hildebrand, and P.~J.~C. Dickinson,
\emph{The extreme rays of the $6\times6$ copositive cone},
J. Global Optim. \textbf{79} (2021), 153--190,
\doi{10.1007/s10898-020-00930-y}.

\bibitem{Dickinson}
P.~J.~C. Dickinson,
\emph{Geometry of the copositive and completely positive cones},
J. Math. Anal. Appl. \textbf{380} (2011), 377--395,
\doi{10.1016/j.jmaa.2011.03.005}.

\bibitem{HA}
R.~Hildebrand and A.~Afonin,
\emph{On the structure of the $6\times6$ copositive cone},
Linear Algebra Appl. \textbf{693} (2024), 22--38,
\doi{10.1016/j.laa.2023.02.004}.

\bibitem{Hildebrand}
R.~Hildebrand,
\emph{Minimal zeros of copositive matrices},
Linear Algebra Appl. \textbf{459} (2014), 154--174,
\doi{10.1016/j.laa.2014.07.004}.

\bibitem{HZ}
T.~Holmgren and Q.~Zhang,
\emph{On dimensions of maximal faces of completely positive cones},
Electron. J. Linear Algebra \textbf{41} (2025), 380--392,
\doi{10.13001/ela.2025.9495}.

\bibitem{KT}
O.~I. Kostyukova and T.~V. Tchemisova,
\emph{Refined estimates on the dimensions of maximal faces of completely
positive cones},
Linear Algebra Appl. (2026), in press,
\doi{10.1016/j.laa.2026.08.010}; arXiv:2603.09633.

\end{thebibliography}

\end{document}